% GENERATED FILE. Edit canonical lessons, then run npm run generate:books.
% canonical-source-hash: fnv1a64:94c5d31ca1401536
% canonical-lessons: BN-C268-mone-hoy, BN-C268-karon, BN-W268-letter-murdhanya-na, BN-C268-karon-jonno, BN-R268-opinion-and-reason

\chapter{What I Think, and Why}
\label{ch:bn-what-i-think-and-why}
\hlchaptermodality{268}

\begin{chapteropening}
\textbf{By the end of this chapter:} \emph{I can say what I think with āmār mône hôy, give a reason with kārôn, kenonā and -er jônyô, and read and write \bn{ণ} in \bn{কারণ}.}

Complete BN-R268-opinion-and-reason: say what you think and give your reason.
\end{chapteropening}

\section[\texorpdfstring{\bn{আমার} \bn{মনে} \bn{হয়} … (āmār mône hôy …)}{আমার মনে হয় … (āmār mône hôy …)}]{\bn{আমার} \bn{মনে} \bn{হয়} … (āmār mône hôy …) — I think}

\label{lesson:BN-C268-mone-hoy}



\begin{quote}
Say \textquotedblleft{}first\textquotedblright{}, \textquotedblleft{}after that\textquotedblright{}, \textquotedblleft{}in the end\textquotedblright{}. (\emph{prôthôme}, \emph{tārpôr}, \emph{sheshe}.) Then \textbf{\bn{মনে} \bn{রাখা}}, \textquotedblleft{}to remember\textquotedblright{}: \textquotedblleft{}to keep in the mind\textquotedblright{}. Its \textbf{\bn{মনে}}, \textquotedblleft{}in the mind\textquotedblright{}, is in this lesson too.
\end{quote}

\subsection*{You'll want to know: I think}
\begin{quote}
\textbf{\bn{আমার} \bn{মনে} \bn{হয়}, \bn{কাল} \bn{ঝড়} \bn{হবে।}}
\end{quote}

\begin{quote}
\emph{āmār mône hôy, kāl jhôṛ hôbe.} — \textquotedblleft{}I think there will be a storm tomorrow.\textquotedblright{}
\end{quote}

Bengali does not say \textquotedblleft{}I think\textquotedblright{} with \textbf{\bn{আমি}}. It says \textquotedblleft{}to my mind it comes\textquotedblright{}: \textbf{\bn{আমার} \bn{মনে} \bn{হয়}}, with the thinker in the \textbf{-\bn{র}} form, as in \textbf{\bn{আমার} \bn{নাম}}. \textbf{\bn{যে}}, \textquotedblleft{}that\textquotedblright{}, which you know, may follow it: \textbf{\bn{আমার} \bn{মনে} \bn{হয়} \bn{যে} …}. To ask, put the other person in the same form: \textbf{\bn{তোমার} \bn{কী} \bn{মনে} \bn{হয়}?} (\emph{tomār kī mône hôy?}), \textquotedblleft{}what do you think?\textquotedblright{}.

\subsection*{You'll want to know: I don't think so}
\textbf{\bn{না}} after \textbf{\bn{হয়}} turns it: \textbf{\bn{আমার} \bn{মনে} \bn{হয়} \bn{না}} (\emph{āmār mône hôy nā}), \textquotedblleft{}I don't think so\textquotedblright{}, and \textbf{\bn{আমার} \bn{মনে} \bn{হয়} \bn{না} \bn{সে} \bn{আসবে}}, \textquotedblleft{}I don't think he'll come\textquotedblright{}. In the past, \textbf{\bn{হয়}} becomes \textbf{\bn{হল}}: \textbf{\bn{আমার} \bn{মনে} \bn{হল}} (\emph{āmār mône hôlô}), \textquotedblleft{}I thought (then), it seemed to me\textquotedblright{}.

\subsection*{Your turn}
Say these aloud:

\begin{itemize}
\raggedright
  \item \textquotedblleft{}I think there will be a storm tomorrow\textquotedblright{} — \emph{āmār mône hôy, kāl jhôṛ hôbe}
  \item \textquotedblleft{}What do you think?\textquotedblright{} — \emph{tomār kī mône hôy?}
  \item \textquotedblleft{}I don't think the train will come\textquotedblright{} — \emph{āmār mône hôy nā ṭren āsbe}
\end{itemize}

\begin{tcolorbox}[breakable,colback=teal!4,colframe=teal!35!black,title={Before you move on}]
\textquotedblleft{}I think\textquotedblright{}? (\textbf{\bn{আমার} \bn{মনে} \bn{হয়}}.) Why \textbf{\bn{আমার}} and not \textbf{\bn{আমি}}? (The thought comes to me.) \textquotedblleft{}I thought, then\textquotedblright{}? (\textbf{\bn{আমার} \bn{মনে} \bn{হল}}.)
\end{tcolorbox}
\section[\texorpdfstring{\bn{কারণ} (kārôn)}{কারণ (kārôn)}]{\bn{কারণ} (kārôn) — a reason; because}

\label{lesson:BN-C268-karon}



\begin{quote}
Before the new one: say \textquotedblleft{}I think\textquotedblright{}. (\emph{āmār mône hôy}.) Then \textbf{\bn{কেননা}}, \textquotedblleft{}because\textquotedblright{}, from chapter twenty-nine.
\end{quote}

\subsection*{You'll want to know: \bn{কারণ}}
\textbf{\bn{কারণ}} — \emph{kārôn} — \textquotedblleft{}a reason; because\textquotedblright{}.

\begin{quote}
\textbf{\bn{তোমার} \bn{রাগের} \bn{কারণ} \bn{কী}?}
\end{quote}

\begin{quote}
\emph{tomār rāger kārôn kī?} — \textquotedblleft{}What's the reason you're angry?\textquotedblright{}
\end{quote}

As a noun it is \textquotedblleft{}a reason\textquotedblright{}: \textbf{\bn{রাগের} \bn{কারণ}}, \textquotedblleft{}the reason for the anger\textquotedblright{}. It is also the everyday \textquotedblleft{}because\textquotedblright{}, opening the reason after the result:

\begin{quote}
\textbf{\bn{আজ} \bn{আমি} \bn{বাড়িতে} \bn{থাকব}, \bn{কারণ} \bn{বাইরে} \bn{ঝড়} \bn{হচ্ছে।}}
\end{quote}

\begin{quote}
\emph{āj āmi bāṛite thākbô, kārôn bāire jhôṛ hôchchhe.} — \textquotedblleft{}I'll stay at home today, because there's a storm outside.\textquotedblright{}
\end{quote}

\begin{cousinweb}
\textbf{\bn{কারণ}} (\emph{kārôn}) is a tatsama: Sanskrit \emph{kāraṇa}, \textquotedblleft{}cause\textquotedblright{}, which CDIAL (3057) lists with the Bengali word under it. CDIAL files \emph{kāraṇa} under the root \emph{kṛ}, \textquotedblleft{}to do, to make\textquotedblright{}: a cause is what makes a thing happen.
\end{cousinweb}

\subsection*{Your turn}
Say these aloud:

\begin{itemize}
\raggedright
  \item \emph{kārôn}
  \item \textquotedblleft{}What's the reason you're angry?\textquotedblright{} — \emph{tomār rāger kārôn kī?}
  \item \textquotedblleft{}I'll stay at home, because there's a storm outside\textquotedblright{} — \emph{āmi bāṛite thākbô, kārôn bāire jhôṛ hôchchhe}
\end{itemize}

\begin{tcolorbox}[breakable,colback=teal!4,colframe=teal!35!black,title={Before you move on}]
What does \textbf{\bn{কারণ}} mean? (\textquotedblleft{}A reason\textquotedblright{}, and \textquotedblleft{}because\textquotedblright{}.) Which comes first, the result or the reason? (The result; \textbf{\bn{কারণ}} opens the reason.)
\end{tcolorbox}
\section[\texorpdfstring{\bn{ণ} (n)}{ণ (n)}]{\bn{ণ} (n) — the letter at the end of \bn{কারণ}}

\label{lesson:BN-W268-letter-murdhanya-na}



\begin{quote}
Before the new one: say \textbf{\bn{কারণ}}, \textquotedblleft{}because\textquotedblright{}. Then read \textbf{\bn{ন}} (\emph{nô}), the first letter you ever wrote.
\end{quote}

\subsection*{Script you'll notice: \bn{ণ}}
\textbf{\bn{ণ}} — \emph{n}. Bengali says \textbf{\bn{ণ}} and \textbf{\bn{ন}} alike, so, like \textbf{\bn{ষ}} beside \textbf{\bn{শ}}, it is a matter of spelling: most words that carry \textbf{\bn{ণ}} come from Sanskrit.

Set the print of \textbf{\bn{ণ}} beside \textbf{\bn{ন}}. Both have an upright stem on the right and an arch on the left that ends in a solid knob. In \textbf{\bn{ন}}, the headline is long and covers the whole letter, and the arch meets the stem low down, near its middle. In \textbf{\bn{ণ}}, the headline is short and sits only at the stem and to its right; the arch meets the stem near the top, where the headline would be, and the stem's top end stands a little above the short headline, cut on a slant.

\subsection*{Writing: \bn{ণ} — copy the printed shape}
Put your pen beside \textbf{\bn{ণ}} and copy the shape you can see, slowly and large. Check the copy against the print: the short headline right of the stem, the arch joined near the top, the knob, the stem to its foot.

\begin{quote}
This book does not yet tell you where to start \bn{ণ} or which way the pen goes. That is not written down here until it can be written down with a source. Copying the printed shape needs no such source, and it is how the shape becomes yours in the meantime.
\end{quote}

\subsection*{Your turn}
\begin{itemize}
\raggedright
  \item \emph{Find:} \bn{ণ} at the end of \textbf{\bn{কারণ}}
  \item \emph{Write it:} \bn{ন}, then \bn{ণ}, from the print, naming each
  \item \emph{Write it:} \textbf{\bn{কারণ}}
  \item \emph{Write it:} \textbf{\bn{হঠাৎ}}, from memory
\end{itemize}

\begin{tcolorbox}[breakable,colback=teal!4,colframe=teal!35!black,title={Before you move on}]
Which letter ends \textbf{\bn{কারণ}}? (\textbf{\bn{ণ}}.) Where does its headline sit? (Only at the stem and to its right.) Where does its arch meet the stem? (Near the top; in \textbf{\bn{ন}}, low down.)
\end{tcolorbox}
\section[\texorpdfstring{… \bn{কারণ} … / -\bn{এর} \bn{জন্য} (… kārôn … / -er …)}{… কারণ … / -এর জন্য (… kārôn … / -er …)}]{… \bn{কারণ} … / -\bn{এর} \bn{জন্য} (kārôn / -er jônyô) — giving the reason}

\label{lesson:BN-C268-karon-jonno}



\begin{quote}
Say \textquotedblleft{}I'll stay at home, because there's a storm outside\textquotedblright{}. (\emph{āmi bāṛite thākbô, kārôn bāire jhôṛ hôchchhe}.)
\end{quote}

\subsection*{You'll want to know: because, with a sentence}
\begin{quote}
\textbf{\bn{আজ} \bn{আমি} \bn{যাব} \bn{না}, \bn{কারণ} \bn{আমি} \bn{ক্লান্ত।}}
\end{quote}

\begin{quote}
\emph{āj āmi jābô nā, kārôn āmi klāntô.} — \textquotedblleft{}I won't go today, because I'm tired.\textquotedblright{}
\end{quote}

A whole reason, with its own verb or its own \textbf{\bn{ক্লান্ত}}, follows \textbf{\bn{কারণ}}. \textbf{\bn{কেননা}}, which you met in chapter twenty-nine, does the same job in the same place: \textbf{\bn{আজ} \bn{আমি} \bn{যাব} \bn{না}, \bn{কেননা} \bn{আমি} \bn{ক্লান্ত}}. To ask for one, \textbf{\bn{কেন}?} (\emph{kenô?}), \textquotedblleft{}why?\textquotedblright{}, the word inside \textbf{\bn{কেননা}}.

\subsection*{You'll want to know: because of, for}
For a reason that is only a noun, Bengali puts the noun in its \textbf{-\bn{র}} form and adds \textbf{\bn{জন্য}} (\emph{jônyô}): \textbf{-\bn{এর} \bn{জন্য}}, \textquotedblleft{}because of\textquotedblright{}.

\begin{quote}
\textbf{\bn{ঝড়ের} \bn{জন্য} \bn{ট্রেন} \bn{দেরিতে} \bn{এল।}}
\end{quote}

\begin{quote}
\emph{jhôṛer jônyô ṭren derite elô.} — \textquotedblleft{}The train came late because of the storm.\textquotedblright{}
\end{quote}

The same \textbf{\bn{জন্য}} is \textquotedblleft{}for\textquotedblright{}: \textbf{\bn{এটা} \bn{তোমার} \bn{জন্য}} (\emph{eṭā tomār jônyô}), \textquotedblleft{}this is for you\textquotedblright{}. After a vowel the ending is \textbf{-\bn{র}}: \textbf{\bn{মীরার} \bn{জন্য}}.

\subsection*{Your turn}
Say these aloud:

\begin{itemize}
\raggedright
  \item \textquotedblleft{}Why? Because I'm tired.\textquotedblright{} — \emph{kenô? kārôn āmi klāntô.}
  \item \textquotedblleft{}The train came late because of the storm\textquotedblright{} — \emph{jhôṛer jônyô ṭren derite elô}
  \item \textquotedblleft{}This is for you\textquotedblright{} — \emph{eṭā tomār jônyô}
\end{itemize}

\begin{tcolorbox}[breakable,colback=teal!4,colframe=teal!35!black,title={Before you move on}]
What follows \textbf{\bn{কারণ}}? (A whole reason.) What follows a noun? (\textbf{-\bn{এর} \bn{জন্য}}.) \textquotedblleft{}Why?\textquotedblright{} (\textbf{\bn{কেন}?})
\end{tcolorbox}
\section[(review)]{Review — what I think, and why}

\label{lesson:BN-R268-opinion-and-reason}



\begin{quote}
Say \textquotedblleft{}I think there will be a storm tomorrow\textquotedblright{}, then \textquotedblleft{}I won't go, because I'm tired\textquotedblright{}. (\emph{āmār mône hôy, kāl jhôṛ hôbe}; \emph{āmi jābô nā, kārôn āmi klāntô}.)
\end{quote}

\subsection*{Your turn}
\begin{itemize}
\raggedright
  \item \emph{Say it:} What do you think? I think Mira is right.
  \item \emph{Say it:} I don't think the bus will come today.
  \item \emph{Say it:} I stayed at home yesterday, because I had a fever.
  \item \emph{Say it:} We came late because of the storm.
  \item \emph{Say it:} What was the reason? I don't know.
  \item \emph{Write it:} \textbf{\bn{কারণ}}, then \textbf{\bn{ঠিক}}
\end{itemize}

\begin{tcolorbox}[breakable,colback=teal!4,colframe=teal!35!black,title={Before you move on}]
\textquotedblleft{}I don't think so\textquotedblright{}? (\textbf{\bn{আমার} \bn{মনে} \bn{হয়} \bn{না}}.) Which word gives a reason with a verb? (\textbf{\bn{কারণ}}, or \textbf{\bn{কেননা}}.) Which gives one with a noun? (\textbf{-\bn{এর} \bn{জন্য}}.) Which letter ends \textbf{\bn{কারণ}}? (\textbf{\bn{ণ}}.)
\end{tcolorbox}
